\chapter{بنية المعالجات الحديثة}

\begin{khulasa}[: ماذا سنتعلم في هذا الفصل؟]
\begin{itemize}
    \item كيف تطورت المعالجات، ولماذا توقف ارتفاع التردد.
    \item المعالجات متعددة الأنوية والمعالجة المتوازية: الفوائد والصعوبات، وقانون \textenglish{Amdahl}.
    \item الذاكرة المخبئية وتسلسل الذاكرة، ولماذا يكون ترتيب قراءة المعطيات مهمًا جدًا.
\end{itemize}
\end{khulasa}

\begin{intabih}
    كل الأرقام في هذا الفصل \textbf{قسناها فعلًا} على حاسوب محمول بمعالج \textenglish{Intel Core i7-8665U}
    (4 أنوية، 8 خيوط، ذاكرة مخبئية \textenglish{L1 = 32 KB} و\textenglish{L2 = 256 KB} لكل نواة، و\textenglish{L3 = 8 MB} مشتركة).
    إذا أعدت التجارب على حاسوبك ستحصل على أرقام مختلفة قليلًا، لكن الظواهر نفسها.
\end{intabih}

%================================================
\section{المعالجات الحديثة: الاتجاهات والتطور}
%================================================

\subsection{خمسون سنة من التقدم}
\begin{center}
\small
\begin{tabular}{|c|c|c|c|c|}
    \hline
    \ras{المعالج} & \ras{السنة} & \ras{الترانزستورات} & \ras{التردد} & \ras{الأنوية} \\ \hline
    \textenglish{Intel 4004} & 1971 & \textenglish{2\,300} & \textenglish{0.74 MHz} & 1 \\ \hline
    \textenglish{Intel 8086} & 1978 & \textenglish{29\,000} & \textenglish{5--10 MHz} & 1 \\ \hline
    \textenglish{Intel 80386} & 1985 & \textenglish{275\,000} & \textenglish{16--33 MHz} & 1 \\ \hline
    \textenglish{Pentium 4} & 2000 & 42 مليون & \textenglish{1.3--3.8 GHz} & 1 \\ \hline
    \textenglish{Core i7-8665U} & 2019 & حوالي ملياري & \textenglish{1.9--4.8 GHz} & 4 \\ \hline
    معالجات 2024 & 2024 & عشرات المليارات & حوالي \textenglish{5 GHz} & من 8 إلى أكثر من 100 \\ \hline
\end{tabular}
\end{center}

\subsection{جدار التردد}
حتى سنة 2004 تقريبًا، كانت كل سنة تأتي بمعالج أسرع لأن \textbf{التردد} يرتفع. ثم توقف التردد عند حوالي
\textenglish{4} إلى \textenglish{5 GHz}. السبب الرئيسي هو \textbf{الحرارة}:
\begin{itemize}
    \item كل نبضة ساعة تستهلك طاقة. كلما ارتفع التردد ارتفعت الطاقة المستهلكة، والطاقة تتحول إلى حرارة.
    \item عند \textenglish{4 GHz} تقريبًا يسخن المعالج إلى درجة لا نستطيع تبريدها في حاسوب عادي.
\end{itemize}

\begin{tashbih}
    عدّاء يركض أسرع وأسرع، حتى يصل إلى سرعة لا يستطيع أن يحافظ عليها دون أن ينهار. بدل أن نطلب منه
    الركض أسرع، نضيف عدّائين آخرين يتقاسمون معه المسافة. هذا بالضبط ما فعله صانعو المعالجات.
\end{tashbih}

\subsection{الاتجاهات الحالية}
\begin{enumerate}
    \item \textbf{أنوية أكثر} بدل تردد أعلى (القسم الآتي).
    \item \textbf{ذاكرة مخبئية أكبر وأذكى}، لأن الذاكرة أصبحت أبطأ بكثير من المعالج (القسم الأخير).
    \item \textbf{تعليمات متخصصة}: \textenglish{SIMD} (الفصل الرابع)، وتسريع الذكاء الاصطناعي والتشفير.
    \item \textbf{أنوية مختلفة في الشريحة نفسها}: أنوية قوية للعمل الثقيل وأنوية اقتصادية للعمل الخفيف، لتوفير البطارية.
\end{enumerate}

%================================================
\section{المعالجات متعددة الأنوية والتوازي}
%================================================

\subsection{المفاهيم}
\begin{tarif}
    \begin{itemize}
        \item \mustalah{النواة}{Cœur}{Core}: معالج كامل (وحدة تحكم، وحدة حساب، سجلات، ذاكرة مخبئية \textenglish{L1}).
        \item \mustalah{المعالج متعدد الأنوية}{Processeur multi-cœur}{Multi-core processor}: عدة أنوية في شريحة واحدة.
        \item \mustalah{الخيط}{Thread}{Thread}: سلسلة من التعليمات تُنفذ بشكل مستقل. البرنامج قد يملك خيطًا واحدًا أو عدة خيوط.
        \item \mustalah{المعالجة المتوازية}{Traitement parallèle}{Parallel processing}: تنفيذ عدة أجزاء من العمل
              \textbf{في اللحظة نفسها} على أنوية مختلفة.
    \end{itemize}
\end{tarif}

\begin{center}
\begin{tikzpicture}[core/.style={unit, minimum width=2.4cm, minimum height=1.5cm}]
    \draw[brandDark, very thick, rounded corners=8pt, fill=brandLight] (-6.1,-2.6) rectangle (6.1,2.1);
    \node[font=\bfseries] at (0,1.75) {شريحة معالج من 4 أنوية};
    \foreach \x/\i in {4.5/1, 1.5/2, -1.5/3, -4.5/4} {
        \node[core] at (\x,0.6) {النواة \i\\{\scriptsize \textenglish{L1}، \textenglish{L2}}};
    }
    \node[unit=brandOrange, minimum width=11.2cm] at (0,-1.1) {الذاكرة المخبئية المشتركة \textenglish{L3}};
    \node[unit=brandPurple, minimum width=4cm] (ram) at (0,-3.6) {الذاكرة الرئيسية \textenglish{RAM}};
    \draw[sahmT, <->] (0,-1.6) -- (ram);
\end{tikzpicture}
\end{center}

\begin{tashbih}
    مطبخ فيه طباخ واحد (نواة واحدة) يحضّر الطبق تلو الآخر. مطبخ فيه 4 طباخين (4 أنوية) يحضّر 4 أطباق
    في الوقت نفسه. لكنهم يتقاسمون الثلاجة نفسها (\textenglish{L3} و\textenglish{RAM})، وقد ينتظر أحدهم دوره
    أمامها. وطبق واحد معقد لا يصبح أسرع إلا إذا استطعنا تقسيمه إلى مهام مستقلة.
\end{tashbih}

\subsection{\textenglish{Hyper-Threading}: نواة واحدة لخيطين}
تقنية \mustalah{تعدد الخيوط المتزامن}{Multithreading simultané}{SMT, Hyper-Threading} تجعل النواة الواحدة تبدو
كأنها نواتان: عندما ينتظر خيط معطيات من الذاكرة، يستعمل الخيط الثاني وحدات الحساب الفارغة. لهذا يُظهر
معالجنا 8 «معالجات منطقية» رغم أنه يملك 4 أنوية فقط. لكن إذا كان الخيطان يحتاجان وحدة الحساب نفسها
دائمًا، فلا نربح شيئًا تقريبًا، كما سنرى في التجربة.

\subsection{تجربة: العمل نفسه على 1 و2 و4 و8 خيوط}
البرنامج الآتي يحسب مجموع الجذور التربيعية لـ400 مليون عدد، ويقسم العمل بين $k$ خيوط:

\cfile{ch5/threads.c}
\screen{ch5/threads.out}

\begin{sharh}
    \s{18--26}{كل خيط يأخذ الأعداد $id$، $id + k$، $id + 2k$\dots فيحسب جزءًا من المجموع لا يتداخل مع الآخرين.}
    \s{34--35}{\code{pthread\_create}: نطلق $k$ خيوط، يعمل كل واحد منها على نواة.}
    \s{36--39}{\code{pthread\_join}: ننتظر انتهاء كل خيط، ثم نجمع النتائج الجزئية.}
    \s{40--42}{التسريع $=$ الزمن بخيط واحد $\div$ الزمن بـ$k$ خيوط.}
\end{sharh}

\begin{center}
\begin{tikzpicture}[x=1.1cm, y=1.1cm]
    \draw[->, thick] (0,0) -- (8.6,0) node[right, font=\small] {الخيوط};
    \draw[->, thick] (0,0) -- (0,4.4) node[above, font=\small] {التسريع};
    \foreach \y in {1,2,3,4} \draw (0.08,\y) -- (-0.08,\y) node[left, font=\scriptsize] {\textenglish{\y}};
    \foreach \x in {1,2,4,8} \draw (\x,0.08) -- (\x,-0.08) node[below, font=\scriptsize] {\textenglish{\x}};
    \draw[brandMuted, dashed, thick] (1,1) -- (4,4) node[above, font=\scriptsize, lawn=brandMuted] {التسريع المثالي};
    \draw[orangeDeep, very thick] (1,1) -- (2,1.98) -- (4,3.02) -- (8,3.05);
    \foreach \x/\y in {1/1, 2/1.98, 4/3.02, 8/3.05} {
        \fill[orangeDeep] (\x,\y) circle (2.5pt);
        \node[font=\codefont\scriptsize, above left] at (\x,\y) {\textenglish{\y}};
    }
    \draw[brandTeal, dashed] (4,0) -- (4,4.2);
    \node[font=\scriptsize, lawn=tealDeep, align=center] at (6,3.9) {4 أنوية فقط:\\الخيوط الإضافية لا تساعد هنا};
\end{tikzpicture}
\end{center}

\begin{khatwa}[: كيف نقرأ هذه النتائج؟]
    \begin{enumerate}
        \item \textbf{خيطان}: التسريع \textenglish{1.98}، أي قريب جدًا من 2. العمل انقسم إلى نصفين مستقلين تمامًا.
        \item \textbf{4 خيوط}: التسريع \textenglish{3.02} وليس 4. عندما تعمل الأنوية الأربعة معًا يسخن المعالج،
              فيخفض تردده تلقائيًا. وجزء من البرنامج (إطلاق الخيوط وجمع النتائج) لا يتوازى.
        \item \textbf{8 خيوط}: التسريع \textenglish{3.05}، تقريبًا مثل 4 خيوط. معالجنا فيه 4 أنوية حقيقية فقط،
              والخيوط الإضافية تتقاسم وحدات الحساب نفسها، فلا يتقدم العمل أسرع.
    \end{enumerate}
\end{khatwa}

\subsection{قانون \textenglish{Amdahl}: حدود التوازي}
معظم البرامج فيها جزء \textbf{لا يمكن توزيعه} (قراءة الملف، عرض النتيجة، انتظار المستعمل\dots).

\begin{tarif}
    إذا كان الجزء القابل للتوازي من البرنامج هو $p$ (بين 0 و1)، ولدينا $n$ نواة، فإن التسريع الأقصى هو:
    \[ S(n) = \frac{1}{(1 - p) + \dfrac{p}{n}} \]
\end{tarif}

\begin{khatwa}[: برنامج 90\% منه قابل للتوازي]
    $p = 0.9$، والجزء التسلسلي $1 - p = 0.1$.
    \begin{enumerate}
        \item بـ4 أنوية: $S = \dfrac{1}{0.1 + \frac{0.9}{4}} = \dfrac{1}{0.1 + 0.225} = \dfrac{1}{0.325} \approx 3.08$.
        \item بـ16 نواة: $S = \dfrac{1}{0.1 + 0.05625} \approx 6.4$.
        \item بعدد لا نهائي من الأنوية: $\dfrac{p}{n} \rightarrow 0$، فيصبح $S = \dfrac{1}{0.1} = 10$.
    \end{enumerate}
    مهما أضفنا أنوية، لن يكون البرنامج أسرع من 10 مرات، لأن الـ10\% التسلسلية تبقى كما هي.
\end{khatwa}

\begin{hal}
    لاحظ أن قانون \textenglish{Amdahl} مع $p = 0.9$ و4 أنوية يعطي \textenglish{3.08}، وهو قريب جدًا من التسريع
    \textenglish{3.02} الذي قسناه! القانون البسيط يشرح الواقع جيدًا.
\end{hal}

\subsection{الفوائد والتحديات}
\begin{center}
\small
\begin{tabularx}{\linewidth}{|R|R|}
    \hline
    \ras{الفوائد} & \ras{التحديات} \\ \hline
    أداء أكبر دون رفع التردد & يجب أن يُكتب البرنامج بخيوط متعددة، وهذا أصعب \\ \hline
    استهلاك طاقة أقل لنفس العمل & الجزء التسلسلي يحد من التسريع (قانون \textenglish{Amdahl}) \\ \hline
    تشغيل عدة برامج في الوقت نفسه بسلاسة & \textbf{التزامن}: خيطان يكتبان في المتغير نفسه قد يفسدان النتيجة \\ \hline
    إذا تعطلت مهمة، تبقى الأنوية الأخرى تعمل & \textbf{تماسك الذاكرة المخبئية}: يجب أن ترى كل الأنوية القيمة نفسها للمتغير \\ \hline
\end{tabularx}
\end{center}

\begin{mithal}[: مشكلة التزامن]
    خيطان ينفذان \code{x = x + 1} على المتغير نفسه، و\textenglish{x = 5} في البداية:
    \begin{enumerate}
        \item الخيط 1 يقرأ \textenglish{x = 5}.
        \item الخيط 2 يقرأ \textenglish{x = 5} أيضًا (قبل أن يكتب الأول).
        \item الخيط 1 يحسب $5 + 1$ ويكتب 6.
        \item الخيط 2 يحسب $5 + 1$ ويكتب 6 أيضًا.
    \end{enumerate}
    النتيجة 6 بدل 7: ضاعت زيادة واحدة! لهذا نستعمل \textbf{الأقفال} \textenglish{(mutex)} أو تعليمات \textbf{ذرّية}
    (\code{lock add} في \textenglish{x86}) تمنع خيطين من تعديل المتغير في الوقت نفسه. في برنامجنا تجنبنا المشكلة:
    كل خيط يكتب في خانته الخاصة \code{part[id]}.
\end{mithal}

%================================================
\section{الذاكرة المخبئية وتسلسل الذاكرة}
%================================================

\subsection{المشكلة: الذاكرة بطيئة جدًا}
المعالج قادر على تنفيذ عدة تعليمات في كل نانوثانية، لكن قراءة الذاكرة الرئيسية تحتاج حوالي 100 نانوثانية.
لو انتظر المعالج الذاكرة في كل مرة، لبقي معظم الوقت دون عمل.

\begin{tarif}
    \mustalah{الذاكرة المخبئية}{Mémoire cache, antémémoire}{Cache memory}: ذاكرة صغيرة وسريعة جدًا (من نوع
    \textenglish{SRAM}) قريبة من المعالج، تحفظ \textbf{نسخة} من المعطيات التي استُعملت مؤخرًا أو التي سنحتاجها قريبًا.
\end{tarif}

\begin{tashbih}
    عندما تراجع للامتحان، لا تضع كل كتب المكتبة على مكتبك. تضع على المكتب الكتب التي تحتاجها الآن (\textenglish{L1})،
    وفي الحقيبة كتب الأسبوع (\textenglish{L2})، وفي خزانة البيت كتب السنة (\textenglish{L3})، والباقي في المكتبة
    البعيدة (\textenglish{RAM}). كلما كان المكان أقرب، كان أصغر وأسرع.
\end{tashbih}

\subsection{تسلسل الذاكرة}
\begin{center}
\begin{tikzpicture}
    \foreach \y/\w/\t/\c/\d in {
        3.6/1.6/{السجلات}/brandOrange/{أقل من \textenglish{1 KB}، أقل من \textenglish{1 ns}},
        2.7/2.8/{\textenglish{L1}}/brandOrange!70!brandTeal/{\textenglish{32 KB}، حوالي \textenglish{1 ns}},
        1.8/4.0/{\textenglish{L2}}/brandTeal/{\textenglish{256 KB}، حوالي \textenglish{4 ns}},
        0.9/5.2/{\textenglish{L3}}/brandTeal!70!brandPurple/{\textenglish{8 MB}، حوالي \textenglish{10 ns}},
        0.0/6.4/{الذاكرة الرئيسية \textenglish{RAM}}/brandPurple/{\textenglish{8 GB}، حوالي \textenglish{100 ns}},
        -0.9/7.6/{القرص \textenglish{SSD}}/brandMuted/{\textenglish{256 GB}، حوالي \textenglish{100 µs}}} {
        \fill[\c!35, draw=\c, thick] (-\w/2,\y) rectangle (\w/2,\y+0.8);
        \node[font=\small\bfseries] at (0,\y+0.4) {\t};
        \node[font=\scriptsize, anchor=east] at (-4.1,\y+0.4) {\d};
    }
    \draw[sahmO] (4.3,-0.9) -- (4.3,4.4) node[midway, right, font=\small, align=left] {أسرع\\أصغر\\أغلى};
\end{tikzpicture}
\end{center}

\subsection{تجربة: نرى طبقات الذاكرة بأعيننا}
البرنامج الآتي يقرأ جدولًا بترتيب عشوائي (كل قراءة تعطي عنوان القراءة التالية، فلا يستطيع المعالج التخمين)،
ويقيس متوسط زمن القراءة الواحدة، لأحجام مختلفة من الجدول:

\cfile{ch5/latency.c}
\screen{ch5/latency.out}

\begin{center}
\begin{tikzpicture}[x=1.45cm, y=0.042cm]
    \draw[->, thick] (0,0) -- (7.4,0) node[right, font=\small] {حجم المعطيات};
    \draw[->, thick] (0,0) -- (0,125) node[above, font=\small] {الزمن \textenglish{(ns)}};
    \foreach \y in {0,25,50,75,100} \draw (0.06,\y) -- (-0.06,\y) node[left, font=\scriptsize] {\textenglish{\y}};
    \foreach \x/\t in {1/16 KB, 2/64 KB, 3/256 KB, 4/1 MB, 5/4 MB, 6/16 MB, 7/64 MB}
        \node[font=\codefont\scriptsize, below] at (\x,-2) {\textenglish{\t}};
    \fill[brandOrange!15] (0.5,0) rectangle (1.6,125); \node[font=\scriptsize] at (1.05,118) {\textenglish{L1}};
    \fill[brandTeal!15] (1.6,0) rectangle (3.4,125); \node[font=\scriptsize] at (2.5,118) {\textenglish{L2}};
    \fill[brandPurple!12] (3.4,0) rectangle (5.6,125); \node[font=\scriptsize] at (4.5,118) {\textenglish{L3}};
    \fill[brandMuted!12] (5.6,0) rectangle (7.3,125); \node[font=\scriptsize] at (6.45,118) {\textenglish{RAM}};
    \draw[orangeDeep, very thick] (1,1.7) -- (2,2.2) -- (3,4.1) -- (4,9.2) -- (5,12.8) -- (6,83.5) -- (7,112.7);
    \foreach \x/\y in {1/1.7, 2/2.2, 3/4.1, 4/9.2, 5/12.8, 6/83.5, 7/112.7} {
        \fill[orangeDeep] (\x,\y) circle (2.5pt);
        \node[font=\codefont\scriptsize, above] at (\x,\y+2) {\textenglish{\y}};
    }
\end{tikzpicture}
\end{center}

\begin{khatwa}[: ماذا نرى في المنحنى؟]
    \begin{enumerate}
        \item \textbf{حتى \textenglish{32 KB}}: الجدول كله في \textenglish{L1}، والقراءة تستغرق \textenglish{1.7 ns} فقط.
        \item \textbf{\textenglish{64 KB} و\textenglish{256 KB}}: الجدول لم يعد يدخل في \textenglish{L1}، فنقرأ من \textenglish{L2}: الزمن يرتفع إلى \textenglish{2} ثم \textenglish{4 ns}.
        \item \textbf{\textenglish{1 MB} و\textenglish{4 MB}}: نقرأ من \textenglish{L3} المشتركة: حوالي \textenglish{9} إلى \textenglish{13 ns}.
        \item \textbf{\textenglish{16 MB} فأكثر}: الجدول أكبر من \textenglish{L3} (\textenglish{8 MB})، فنذهب إلى \textenglish{RAM}:
              الزمن يقفز إلى \textenglish{84} ثم \textenglish{113 ns}، أي أبطأ بحوالي \textbf{65 مرة} من \textenglish{L1}!
    \end{enumerate}
    الدرجات في المنحنى تقع بالضبط عند أحجام الذاكرات المخبئية لمعالجنا.
\end{khatwa}

\subsection{مبدأ المحلية: لماذا تنجح الذاكرة المخبئية؟}
الذاكرة المخبئية صغيرة، لكنها تحتوي في الغالب على ما نحتاجه، بفضل سلوكين طبيعيين في البرامج:
\begin{itemize}
    \item \mustalah{المحلية الزمنية}{Localité temporelle}{Temporal locality}: المعطاة التي استعملناها الآن سنستعملها
          غالبًا \textbf{مرة أخرى قريبًا} (عدّاد الحلقة، المجموع\dots).
    \item \mustalah{المحلية المكانية}{Localité spatiale}{Spatial locality}: بعد قراءة عنوان، نقرأ غالبًا
          \textbf{العناوين المجاورة} (عناصر جدول، تعليمات متتالية).
\end{itemize}

\begin{tarif}
    \mustalah{سطر الذاكرة المخبئية}{Ligne de cache}{Cache line}: الذاكرة المخبئية لا تنقل بايتًا واحدًا، بل
    كتلة كاملة من \textbf{64 بايتًا} مجاورة. عندما نقرأ \code{t[0]}، تصل معه \code{t[1]} إلى \code{t[15]} (للأعداد من 4 بايتات) مجانًا.
\end{tarif}

\subsection{تجربة: ترتيب القراءة يغيّر كل شيء}
\cfile{ch5/locality.c}
\screen{ch5/locality.out}

\begin{center}
\begin{tikzpicture}[c/.style={draw=brandDark, minimum size=0.55cm}]
    \foreach \i in {0,...,3} \foreach \j in {0,...,3} \node[c] at (-\j*0.55-4.2, -\i*0.55) {};
    \foreach \i in {0,...,3} \draw[sahmT] (-6.0, -\i*0.55) -- (-4.0, -\i*0.55);
    \node[font=\small, align=center] at (-5.0,0.75) {سطرًا بسطر: كل قراءة\\بجوار التي قبلها};
    \foreach \i in {0,...,3} \foreach \j in {0,...,3} \node[c] at (\j*0.55+3.0, -\i*0.55) {};
    \foreach \j in {0,...,3} \draw[sahmO] (\j*0.55+3.0, 0.2) -- (\j*0.55+3.0, -1.85);
    \node[font=\small, align=center] at (3.8,0.75) {عمودًا بعمود: كل قراءة\\تقفز \textenglish{16 KB}};
\end{tikzpicture}
\end{center}

\begin{khatwa}[: لماذا الأعمدة أبطأ بـ33 مرة؟]
    \begin{enumerate}
        \item في لغة \textenglish{C}، الجدول \code{m[4096][4096]} محفوظ \textbf{سطرًا بعد سطر}: \code{m[0][0]} ثم \code{m[0][1]} ثم\dots
        \item \textbf{سطرًا بسطر}: كل سطر ذاكرة مخبئية (64 بايتًا) يحتوي على 16 عددًا متجاورًا. نقرأ من \textenglish{RAM} مرة، ثم نستعمل
              16 عددًا من الذاكرة المخبئية. والمعالج يلاحظ أننا نتقدم بانتظام فيجلب الأسطر القادمة مسبقًا.
        \item \textbf{عمودًا بعمود}: \code{m[0][j]} و\code{m[1][j]} بينهما $4096 \times 4 = 16\,384$ بايتًا. كل قراءة
              في سطر ذاكرة مخبئية جديد، فنستعمل عددًا واحدًا فقط من كل 16 عددًا نقلناها.
        \item النتيجة الحقيقية: \textenglish{0.009 s} مقابل \textenglish{0.301 s}، أي حوالي \textbf{33 مرة} أبطأ،
              لنفس العمليات تمامًا!
    \end{enumerate}
\end{khatwa}

\subsection{الإصابة والإخفاق وزمن الوصول المتوسط}
\begin{tarif}
    \begin{itemize}
        \item \mustalah{الإصابة}{Succès}{Hit}: المعطاة موجودة في الذاكرة المخبئية.
        \item \mustalah{الإخفاق}{Défaut}{Miss}: المعطاة غير موجودة، فنذهب إلى المستوى التالي (الأبطأ).
        \item \textbf{معدل الإصابة} $h$: نسبة القراءات التي وجدت معطياتها في الذاكرة المخبئية.
    \end{itemize}
    \[ \text{الزمن المتوسط} = h \times t_{cache} + (1 - h) \times t_{RAM} \]
\end{tarif}

\begin{khatwa}[: \textenglish{$t_{cache} = 1$ ns}، \textenglish{$t_{RAM} = 100$ ns}]
    \begin{enumerate}
        \item معدل إصابة 90\%: $0.9 \times 1 + 0.1 \times 100 = 0.9 + 10 = 10.9$ \textenglish{ns}.
        \item معدل إصابة 99\%: $0.99 \times 1 + 0.01 \times 100 = 0.99 + 1 = 1.99$ \textenglish{ns}.
    \end{enumerate}
    الانتقال من 90\% إلى 99\% يجعل الذاكرة \textbf{أسرع بأكثر من 5 مرات}. لهذا تهتم المعالجات الحديثة جدًا
    بتحسين معدل الإصابة.
\end{khatwa}

\subsection{أين نضع المعطاة في الذاكرة المخبئية؟ (التطابق)}
الذاكرة المخبئية مقسمة إلى أسطر. علينا أن نقرر في أي سطر نضع كتلة قادمة من الذاكرة:

\begin{center}
\small
\begin{tabularx}{\linewidth}{|>{\centering\arraybackslash}p{3.4cm}|R|R|}
    \hline
    \ras{النوع} & \ras{القاعدة} & \ras{الميزة والعيب} \\ \hline
    التطابق المباشر \textenglish{(direct mapped)} & كل كتلة لها سطر واحد ممكن: رقم الكتلة \textenglish{mod} عدد الأسطر & بسيط وسريع، لكن كتلتان قد تتصارعان على السطر نفسه \\ \hline
    التجميعي الكامل \textenglish{(fully associative)} & أي كتلة في أي سطر & أقل صراعات، لكن البحث في كل الأسطر مكلف \\ \hline
    التجميعي بمجموعات \textenglish{(set associative)} & الكتلة تذهب إلى مجموعة محددة، وفيها إلى أي سطر من $k$ أسطر & الحل الوسط المستعمل في كل المعالجات الحديثة \\ \hline
\end{tabularx}
\end{center}

\begin{khatwa}[: تقسيم عنوان في ذاكرة مخبئية بتطابق مباشر]
    ذاكرة مخبئية سعتها \textenglish{1 KB}، وأسطرها من 16 بايتًا، والعناوين من 16 بتًا.
    \begin{enumerate}
        \item عدد الأسطر: $1024 \div 16 = 64 = 2^6$ سطرًا، إذن نحتاج \textbf{6 بتات} لرقم السطر.
        \item داخل السطر 16 بايتًا $= 2^4$، إذن \textbf{4 بتات} للإزاحة.
        \item الباقي: $16 - 6 - 4 = 6$ بتات، وتسمى \textbf{الوسم} \textenglish{(tag)}: تحدد أي كتلة من الذاكرة موجودة الآن في السطر.
    \end{enumerate}
    \begin{center}
    \begin{tikzpicture}[f/.style={draw=brandDark, thick, minimum height=0.7cm, font=\small\bfseries}]
        % drawn left to right, like the binary address: tag first
        \node[f, minimum width=3cm, fill=brandOrange!20] (t) at (0,0) {الوسم: 6 بتات};
        \node[f, minimum width=3cm, fill=brandTeal!20, anchor=west] (l) at (t.east) {السطر: 6 بتات};
        \node[f, minimum width=2cm, fill=brandGreen!20, anchor=west] (o) at (l.east) {الإزاحة: 4 بتات};
        \node[font=\codefont\scriptsize, below=1pt of t.south west, anchor=north west] {\textenglish{bit 15}};
        \node[font=\codefont\scriptsize, below=1pt of o.south east, anchor=north east] {\textenglish{bit 0}};
    \end{tikzpicture}
    \end{center}
    مثال: $\hex{1A2C} = \bin{000110\,100010\,1100}$. الوسم $\bin{000110} = 6$، والسطر $\bin{100010} = 34$، والإزاحة $\bin{1100} = 12$.
\end{khatwa}

\begin{intabih}
    العنوان الثنائي يُكتب دائمًا من اليسار إلى اليمين، حتى في كتاب بالعربية: البتات العليا (الوسم) على اليسار
    والبت 0 على اليمين، تمامًا كما رسمناه في الشكل.
\end{intabih}

\subsection{الاستبدال والكتابة}
\begin{itemize}
    \item \textbf{سياسة الاستبدال}: عندما تمتلئ المجموعة، أي سطر نطرد؟ غالبًا \textenglish{LRU} \textenglish{(Least Recently Used)}:
          السطر الذي لم نستعمله منذ أطول مدة.
    \item \textbf{سياسة الكتابة}:
    \begin{itemize}
        \item \textenglish{Write-through}: نكتب في الذاكرة المخبئية وفي \textenglish{RAM} معًا. بسيط لكنه بطيء.
        \item \textenglish{Write-back}: نكتب في الذاكرة المخبئية فقط، ولا نكتب في \textenglish{RAM} إلا عندما نطرد السطر. أسرع، وهو المستعمل اليوم.
    \end{itemize}
\end{itemize}

\begin{hal}
    تذكر نموذج هارفارد من الفصل الأول: ذاكرة \textenglish{L1} مقسمة إلى جزأين، \textenglish{L1i} للتعليمات
    و\textenglish{L1d} للمعطيات (في معالجنا \textenglish{32 KB} لكل واحدة)، فيستطيع المعالج جلب تعليمة وقراءة معطاة
    في الوقت نفسه. أما \textenglish{L2} و\textenglish{L3} و\textenglish{RAM} فمشتركة بينهما، كما في نموذج فون نيومان.
\end{hal}

%================================================
\section{تمارين}
%================================================

\begin{tamrin}
    برنامج 75\% منه قابل للتوازي. احسب التسريع بـ2 و4 و8 أنوية، والتسريع الأقصى.
\end{tamrin}
\begin{hall}
    $1 - p = 0.25$. \quad $S(2) = \dfrac{1}{0.25 + 0.375} = 1.6$، \quad $S(4) = \dfrac{1}{0.25 + 0.1875} \approx 2.29$،
    \quad $S(8) = \dfrac{1}{0.25 + 0.09375} \approx 2.91$. الأقصى $= \dfrac{1}{0.25} = 4$.
\end{hall}

\begin{tamrin}
    ذاكرة مخبئية زمنها \textenglish{2 ns}، والذاكرة الرئيسية \textenglish{80 ns}، ومعدل الإصابة 95\%. ما الزمن المتوسط؟
\end{tamrin}
\begin{hall}
    $0.95 \times 2 + 0.05 \times 80 = 1.9 + 4 = 5.9$ \textenglish{ns}.
\end{hall}

\begin{tamrin}
    ذاكرة مخبئية بتطابق مباشر: \textenglish{4 KB}، أسطر من 32 بايتًا، عناوين من 32 بتًا. كم بتًا للإزاحة وللسطر وللوسم؟
\end{tamrin}
\begin{hall}
    الإزاحة: $32 = 2^5$، إذن 5 بتات. الأسطر: $4096 \div 32 = 128 = 2^7$، إذن 7 بتات. الوسم: $32 - 7 - 5 = 20$ بتًا.
\end{hall}

\begin{tamrin}
    اشرح لماذا يكون هذا الكود بطيئًا، وكيف نصلحه:
    \begin{codebox}\begin{english}\codefont\small
    for (j = 0; j < N; j++) \\
    \hspace*{2em}for (i = 0; i < N; i++) \\
    \hspace*{4em}s += m[i][j];
    \end{english}\end{codebox}
\end{tamrin}
\begin{hall}
    يقرأ الجدول عمودًا بعمود، فكل قراءة في سطر ذاكرة مخبئية مختلف (لا محلية مكانية). نصلحه بتبديل الحلقتين:
    \code{for (i ...) for (j ...)}، فنقرأ سطرًا بسطر.
\end{hall}

\begin{tamrin}
    لماذا لم يتحسن البرنامج \code{threads.c} عند الانتقال من 4 إلى 8 خيوط في معالجنا؟
\end{tamrin}
\begin{hall}
    لأن المعالج يملك 4 أنوية حقيقية فقط. الخيوط 5 إلى 8 تعمل بتقنية \textenglish{Hyper-Threading} على الأنوية نفسها،
    والبرنامج يستعمل وحدة الحساب طوال الوقت، فلا توجد وحدات فارغة يستفيد منها الخيط الثاني.
\end{hall}
